\chapter*{INTRODUÇÃO}
\addcontentsline{toc}{chapter}{INTRODUÇÃO}

A consolidação da Indústria 4.0 transformou fundamentalmente os paradigmas da manufatura moderna, promovendo a digitalização e a interconexão de sistemas cibernéticos e físicos no chão de fábrica \cite{kagermann2013}. Nesse cenário de constante evolução tecnológica, surge a necessidade de soluções capazes de representar digitalmente os ativos industriais de forma padronizada, interoperável e altamente reutilizável. Para suprir essa demanda semântica, o \textit{Asset Administration Shell} (AAS) foi estabelecido como o padrão central para a implementação de Gêmeos Digitais, permitindo agrupar todas as propriedades, parâmetros e documentações de um ativo sob uma estrutura unificada \cite{zhang2025, olaya2025}.

Apesar da comprovada relevância do modelo AAS para atingir a verdadeira interoperabilidade exigida pela arquitetura RAMI 4.0 \cite{plattform2016}, a sua adoção prática enfrenta barreiras estruturais significativas nas camadas inferiores da pirâmide de automação. As bibliotecas e implementações do AAS atualmente disponíveis no mercado e na academia são, em sua grande maioria, voltadas para ambientes computacionais de alto desempenho, que contam com farta disponibilidade de memória RAM, grande poder de processamento e robusta infraestrutura de sistemas operacionais tradicionais \cite{pribis2021}, desta forma, a adoção de AAS em larga escala é inviável. 

Diante desse contexto, o presente trabalho se propõe a investigar e desenvolver uma solução otimizada de um AAS embarcado. O objetivo é fornecer uma implementação nativa e eficiente de modo a viabilizar, na prática, a adoção do conceito de Gêmeos Digitais em larga escala. Para isso, o documento descreve a construção de uma biblioteca baseada na arquitetura ESP32 com o sistema operacional de tempo real \textit{FreeRTOS}, aliando as capacidades de instrumentação física aos protocolos de comunicação semântica, como o OPC UA \cite{opcfoundation2026}.